% Provenance (never printed): NSI premiere "Algorithmique: parcours sequentiel, tris par insertion et par
% selection, recherche dichotomique; terminaison (variant) et correction (invariant)". See sources/book2/BRIEFS.md.
\chapter{Proving Algorithms Correct}\label{ch:g11:proving-algorithms-correct}

A binary search passes every one of its ten tests. The next day, asked for a number that is not in the list, it never
answers: the program runs until someone stops it. The tests were not wrong; they were too few, as tests always are. A
test checks one input. To be sure of all the inputs, one needs an argument.

This chapter gives the two arguments that settle most loops. The first shows that a loop stops; the second, that what
it computes is right. Both are short, both can be checked in the code itself with \code{assert}, and together they find
the bug of the binary search above, and others like it, before any test does.

\begin{recall}
A \code{while} loop repeats its body as long as its condition is true (Book 1, grade 6). Linear search looks at each
element in turn; binary search halves a sorted list at each step; selection sort and insertion sort put a list in
order (Book 1, grade 8). A specification states what a function needs and what it promises, and an assertion
stops the program if a condition is false (chapter \ref{ch:g10:specifications-and-tests}). Euclid's algorithm finds
the greatest common divisor of two numbers by repeated remainders (mathematics, grade 9).
\end{recall}

\section{Does it stop?}\label{sec:g11:proving-algorithms-correct:termination}

A \code{for} loop over a list stops: it has as many turns as the list has elements. A \code{while} loop stops only if
its condition becomes false, and nothing guarantees that it will.

\begin{definition}[Variant, termination]\label{def:g11:proving-algorithms-correct:variant}
An algorithm has the property of \emph{termination}\index{termination} if it stops on every input that meets its
\omterm{def:g10:specifications-and-tests:pre-post}{precondition}. A \emph{variant}\index{variant (loop variant)} of a loop is a whole number, computed from the
variables, that is never negative while the loop runs and that strictly decreases at each turn.
\end{definition}

A loop with a \omterm{def:g11:proving-algorithms-correct:variant}{variant} stops: a whole number that is at least $0$ and loses at least $1$ at each turn cannot do so more
times than its first value. Euclid's algorithm is an example. Its \omterm{def:g11:proving-algorithms-correct:variant}{variant} is $b$: the new $b$ is the remainder
\texttt{a \% b}, which is smaller than $b$ and never negative.

\omcode{code/grade-11/05-proving-algorithms-correct/listings/variant_check.py}{1}{13}{Euclid's algorithm, printing its
\omterm{def:g11:proving-algorithms-correct:variant}{variant} at each turn.}\label{lst:g11:proving-algorithms-correct:variant}

\omoutput{code/grade-11/05-proving-algorithms-correct/outputs/variant_check.txt}{The \omterm{def:g11:proving-algorithms-correct:variant}{variant} $b$ falls at each turn and
the loop stops when it reaches $0$: three turns, then eleven (the last line of each run is printed after the
loop).}\label{out:g11:proving-algorithms-correct:variant}

\begin{omfigure}
\begin{tikzpicture}
\begin{axis}[omrecord, width=10cm, height=5cm, xlabel={turn}, ylabel={variant $b$}, xmin=-0.5, xmax=12, ymin=0,
    ymax=160, xtick={0,2,...,12}]
  \addplot[const plot, very thick, omDef, mark=*, mark size=1.5pt] coordinates {(0,144) (1,89) (2,55) (3,34) (4,21)
    (5,13) (6,8) (7,5) (8,3) (9,2) (10,1) (11,0)};
\end{axis}
\end{tikzpicture}

{\small The \omterm{def:g11:proving-algorithms-correct:variant}{variant} of \code{gcd(233, 144)}, from the run above: a staircase that only goes down, one step per turn,
and cannot go below $0$.}
\end{omfigure}

The \omterm{def:g11:proving-algorithms-correct:variant}{variant} does not have to fall by $1$ each turn; it may fall by much more. What matters is that it falls by at
least $1$ and starts from a finite value, which bounds the number of turns.

\section{Is the answer right?}\label{sec:g11:proving-algorithms-correct:invariant}

\begin{definition}[Correct algorithm]\label{def:g11:proving-algorithms-correct:correct}
An algorithm is \emph{correct}\index{correct algorithm} with respect to a specification if, on every input that meets
the \omterm{def:g10:specifications-and-tests:pre-post}{precondition}, it stops and its result meets the \omterm{def:g10:specifications-and-tests:pre-post}{postcondition}.
\end{definition}

\begin{definition}[Loop invariant]\label{def:g11:proving-algorithms-correct:invariant}
A \emph{loop invariant}\index{loop invariant} is a statement about the variables that is true before the first
turn of a loop and that each turn keeps true: if it holds at the start of a turn, it holds at the end.
\end{definition}

An invariant is then true when the loop stops, together with the negation of the loop's condition. A good invariant is
chosen so that these two facts, put together, give the \omterm{def:g10:specifications-and-tests:pre-post}{postcondition}.

\begin{method}[Proving a loop with an invariant]
\begin{enumerate}
  \item Write the invariant, often ``the part already handled is right'', with the loop's index.
  \item \textbf{Start}: show that it is true before the first turn.
  \item \textbf{Keep}: suppose it is true at the start of a turn; show that it is still true after the turn.
  \item \textbf{End}: add the fact that the loop's condition is false; deduce the \omterm{def:g10:specifications-and-tests:pre-post}{postcondition}.
  \item Give a \omterm{def:g11:proving-algorithms-correct:variant}{variant}, to show that the end is reached.
\end{enumerate}
\end{method}

The function \code{largest} returns the largest element of a list. Its invariant is ``\code{best} is the largest of
\code{t[:i]}''. Start: with \code{i = 1}, \code{t[:1]} holds only \code{t[0]}, which is \code{best}. Keep: if
\code{t[i]} is larger than \code{best}, it becomes \code{best}, otherwise \code{best} stays; either way it is the
largest of \code{t[:i + 1]}, and \code{i} grows by one. End: the loop stops when \code{i} equals \code{len(t)}, so
\code{t[:i]} is the whole list. \omterm{def:g11:proving-algorithms-correct:variant}{Variant}: \code{len(t) - i}.

\omcode{code/grade-11/05-proving-algorithms-correct/listings/largest.py}{1}{15}{The largest element, with its invariant
written as an assertion.}\label{lst:g11:proving-algorithms-correct:largest}

\omoutput{code/grade-11/05-proving-algorithms-correct/outputs/largest.txt}{No assertion fails on these two lists.}\label{out:g11:proving-algorithms-correct:largest}

An assertion of the invariant does not prove it; it checks it on the inputs that are run, and stops the program at the
first turn where it fails, which points at the bug. The proof covers all the inputs; the assertion catches the gap
between the proof and the code.

\section{Selection sort}\label{sec:g11:proving-algorithms-correct:selection}

Selection sort finds a smallest element of the part not yet sorted and swaps it to the front of that part. Its
invariant has two halves: \code{t[:i]} is sorted, and no element of \code{t[:i]} is larger than an element of
\code{t[i:]}.

\omcode{code/grade-11/05-proving-algorithms-correct/listings/selection_sort.py}{1}{19}{Selection sort, checking both
halves of its invariant.}\label{lst:g11:proving-algorithms-correct:selection}

\omoutput{code/grade-11/05-proving-algorithms-correct/outputs/selection_sort.txt}{The bar separates \code{t[:i + 1]},
sorted after the turn, from the rest.}\label{out:g11:proving-algorithms-correct:selection}

\begin{omfigure}
\begin{tikzpicture}[c/.style={draw, minimum size=8mm, font=\small}]
  \foreach \v [count=\k from 0] in {3, 12, 47, 29, 18, 12} {
    \ifnum\k<2 \def\fillc{omMeth!20}\else\def\fillc{white}\fi
    \node[c, fill=\fillc] (c\k) at (0.8*\k, 0) {\v};
  }
  \node[c, fill=omDef!20] at (c5) {12};
  \draw[decorate, decoration={brace, amplitude=4pt}] (-0.4,0.5) -- node[above=4pt, font=\footnotesize, align=center]
    {\code{t[:2]}\\sorted} (1.2,0.5);
  \draw[decorate, decoration={brace, amplitude=4pt}] (1.2,0.5) -- node[above=4pt, font=\footnotesize, align=center]
    {\code{t[2:]}\\none smaller than $12$} (4.4,0.5);
  \draw[<->, thick, omThm] (c2.south) to[bend right=30] node[below, font=\footnotesize] {swap} (c5.south);
  \node[font=\footnotesize, right] at (4.6,0) {smallest of \code{t[2:]}};
\end{tikzpicture}

{\small Selection sort at the start of turn $2$: the sorted part, the rest, and the swap that will put a smallest
element of the rest in place $2$.}
\end{omfigure}

Start: with \code{i = 0}, \code{t[:0]} is empty, and both halves hold. Keep: the turn puts in place \code{i} a smallest
element of \code{t[i:]}. By the second half, it is at least as large as everything in \code{t[:i]}, so the
first \code{i + 1} elements are sorted; being a smallest of \code{t[i:]}, it is no larger than anything left in \code{t[i + 1:]}. End: with
\code{i = len(t)}, the whole list is sorted. The swap only exchanges two elements, so the list keeps the same elements.
\omterm{def:g11:proving-algorithms-correct:variant}{Variant} of the outer loop: \code{len(t) - i}; the inner loop is a \code{for} loop and stops.

\section{Insertion sort}\label{sec:g11:proving-algorithms-correct:insertion}

Insertion sort takes the elements one by one and slides each into place among those before it, which are already
sorted. It has two loops, and each has its own invariant.

\omcode{code/grade-11/05-proving-algorithms-correct/listings/insertion_sort.py}{1}{14}{Insertion sort, with its outer
invariant asserted.}\label{lst:g11:proving-algorithms-correct:insertion}

\omoutput{code/grade-11/05-proving-algorithms-correct/outputs/insertion_sort.txt}{After turn $i$, the first $i + 1$
elements are sorted.}\label{out:g11:proving-algorithms-correct:insertion}

\begin{omfigure}
\begin{tikzpicture}[c/.style={draw, minimum size=8mm, font=\small}]
  \foreach \v [count=\k from 0] in {3, 12, 29, 47, \phantom{0}, 12} {
    \ifnum\k<4 \def\fillc{omMeth!20}\else\def\fillc{white}\fi
    \node[c, fill=\fillc] (c\k) at (0.8*\k, 0) {\v};
  }
  \node[c, fill=omDef!20] (x) at (3.2, 1.5) {18};
  \draw[->, thick, omThm] (c3.south) to[bend right=60] node[below, font=\scriptsize] {1} (c4.south);
  \draw[->, thick, omThm] (c2.south) to[bend right=60] node[below, font=\scriptsize] {2} (c3.south);
  \draw[->, thick, omDef] (x.west) to[bend right=30] node[above left, font=\footnotesize] {lands in place $2$}
    (c2.north);
  \node[font=\footnotesize, right] at (3.7,1.5) {$x$, taken out of place $4$};
\end{tikzpicture}

{\small Turn $4$ of insertion sort. The sorted part was $3, 12, 29, 47$; the elements larger than $x = 18$ move one place
right, first $47$ then $29$, and $x$ lands in the gap.}
\end{omfigure}

The outer invariant is ``\code{t[:i]} is sorted''. The inner invariant, for the \code{while} loop, is: the elements of
\code{t[j + 1:i + 1]} are larger than $x$, in order, and putting $x$ in place \code{j} would give the old
\code{t[:i]} with $x$ added. Each inner turn moves one larger element right and lowers \code{j}; the inner loop stops
when \code{j} is $0$ or when \code{t[j - 1]} is not larger than $x$. In both cases $x$, written in place \code{j}, is
not smaller than what comes before it and is smaller than what comes after: \code{t[:i + 1]} is sorted, and the outer
invariant is kept. The inner \omterm{def:g11:proving-algorithms-correct:variant}{variant} is \code{j}, the outer one \code{len(t) - i}.

\section{Binary search}\label{sec:g11:proving-algorithms-correct:binary-search}

Binary search finds \code{x} in a sorted list by keeping a window \code{t[lo:hi]} outside of which \code{x} cannot be,
and halving it.

\omcode{code/grade-11/05-proving-algorithms-correct/listings/binary_search.py}{1}{20}{Binary search, its invariant and its
\omterm{def:g11:proving-algorithms-correct:variant}{variant}.}\label{lst:g11:proving-algorithms-correct:binary-search}

\omoutput{code/grade-11/05-proving-algorithms-correct/outputs/binary_search.txt}{Three searches. Looking for $20$, the
window shrinks to nothing and the answer is \code{None}.}\label{out:g11:proving-algorithms-correct:binary-search}

\begin{omfigure}
\begin{tikzpicture}[c/.style={draw, minimum width=7mm, minimum height=6mm, font=\footnotesize}]
  \foreach \row/\lo/\md/\hi in {0/0/5/10, 1/6/8/10, 2/6/7/8} {
    \foreach \v [count=\k from 0] in {3, 8, 12, 15, 21, 30, 34, 41, 47, 52} {
      \pgfmathtruncatemacro{\kind}{(\k < \lo || \k >= \hi) ? 0 : (\k == \md ? 1 : 2)}
      \ifcase\kind \def\fillc{black!12}\or\def\fillc{omDef!25}\else\def\fillc{white}\fi
      \node[c, fill=\fillc] at (0.7*\k, -1.1*\row) {\v};
    }
    \node[font=\footnotesize, anchor=west] at (7.2, -1.1*\row) {\texttt{lo=\lo}, \texttt{mid=\md}, \texttt{hi=\hi}};
  }
\end{tikzpicture}

{\small Looking for $41$: the grey cells are outside the window, where $41$ cannot be; the blue cell is \code{mid}.
The third \code{mid} holds $41$.}
\end{omfigure}

The invariant: \code{x} is neither in \code{t[:lo]} nor in \code{t[hi:]}. Start: both are empty. Keep: if
\code{t[mid] < x}, then, the list being sorted, every element of \code{t[:mid + 1]} is less than \code{x}, so
\code{lo = mid + 1} keeps the invariant; if \code{t[mid] > x}, every element of \code{t[mid:]} is larger, and
\code{hi = mid} keeps it. End: if the loop stops with \code{lo >= hi}, the window is empty, and by the invariant
\code{x} is nowhere: \code{None} is right. The \omterm{def:g11:proving-algorithms-correct:variant}{variant} is \code{hi - lo}: \code{mid} lies in the window, so both
\code{lo = mid + 1} and \code{hi = mid} make it smaller.

The \omterm{def:g11:proving-algorithms-correct:variant}{variant} is where the opening bug hides. Writing \code{lo = mid} keeps the invariant, but not the \omterm{def:g11:proving-algorithms-correct:variant}{variant}: when
\code{hi = lo + 1}, \code{mid} equals \code{lo}, and \code{lo = mid} changes nothing.

\omcode{code/grade-11/05-proving-algorithms-correct/listings/binary_search_bug.py}{1}{25}{The opening bug, with a guard
added for the demonstration.}\label{lst:g11:proving-algorithms-correct:bug}

\omoutput{code/grade-11/05-proving-algorithms-correct/outputs/binary_search_bug.txt}{Every test passes, since each one
searches a number of the list; a number not in it makes the window stop shrinking at \code{lo=3},
\code{hi=4}.}\label{out:g11:proving-algorithms-correct:bug}

\begin{history}[Meaning given to programs]
The idea of attaching statements to the points of a program and proving that each step keeps them true was published
by Robert W. Floyd in 1967, for flowcharts. In 1969, C.~A.~R.~Hoare turned it into rules for reasoning about the
instructions of a programming language, one rule per kind of instruction, the loop rule resting on an invariant.
% ledger: g11:floyd-1967
% ledger: g11:hoare-1969
\end{history}

\section{Exercises}\label{sec:g11:proving-algorithms-correct:exercises}

\begin{exercise}[$\star$]\label{exo:g11:proving-algorithms-correct:1}
Write the list \code{[5, 2, 9, 1, 7]} after the first two turns of selection sort.
\end{exercise}

\begin{exercise}[$\star$]\label{exo:g11:proving-algorithms-correct:2}
Give a \omterm{def:g11:proving-algorithms-correct:variant}{variant} of the loop \code{while n > 0: n = n - 2}, for \code{n} a whole number, and say why it is one.
\end{exercise}

\begin{exercise}[$\star$]\label{exo:g11:proving-algorithms-correct:3}
Give the values of \code{lo}, \code{mid} and \code{hi} at each turn when the chapter's \code{search} looks for $52$
in the list of the output.
\end{exercise}

\begin{exercise}[$\star$]\label{exo:g11:proving-algorithms-correct:4}
For the loop \code{s = 0}, then \code{for i in range(len(t)): s = s + t[i]}, is the statement ``\code{s} is the sum
of \code{t[:i]}'' true before the first turn?
\end{exercise}

\begin{exercise}[$\star$]\label{exo:g11:proving-algorithms-correct:5}
A student writes \code{largest} with \code{for i in range(1, len(t) - 1)}. On \code{[4, 9, 2, 11]} it returns $9$.
Which step of the proof fails?
\end{exercise}

\begin{exercise}[$\star\star$]\label{exo:g11:proving-algorithms-correct:6}
Prove that this loop stops for every whole number \code{n}, by giving a \omterm{def:g11:proving-algorithms-correct:variant}{variant}: \code{while n > 1}, if \code{n} is
even, \code{n = n // 2}, otherwise \code{n = n - 1}.
\end{exercise}

\begin{exercise}[$\star\star$]\label{exo:g11:proving-algorithms-correct:7}
Write \code{index_of_min(t, i)} that returns the index of a smallest element of \code{t[i:]}, ending with an
\code{assert} of its \omterm{def:g10:specifications-and-tests:pre-post}{postcondition}.
\end{exercise}

\begin{exercise}[$\star\star$]\label{exo:g11:proving-algorithms-correct:8}
In the chapter's \code{search}, a student writes \code{hi = mid - 1}. The loop always stops, but on the list of the
output it misses $8$, $21$ and $41$. Which part of the proof fails, and why?
\end{exercise}

\begin{exercise}[$\star\star$]\label{exo:g11:proving-algorithms-correct:9}
Prove, with an invariant, that the loop of exercise \ref{exo:g11:proving-algorithms-correct:4} computes the sum of
\code{t}.
\end{exercise}

\begin{exercise}[$\star\star$]\label{exo:g11:proving-algorithms-correct:10}
Write a loop that counts the occurrences of \code{x} in \code{t}, and state the invariant that proves it.
\end{exercise}

\begin{exercise}[$\star\star$]\label{exo:g11:proving-algorithms-correct:11}
Write \code{is_sorted(t)} with one loop that compares neighbours, and prove it \omterm{def:g11:proving-algorithms-correct:correct}{correct}: it returns \code{True} exactly
when \code{t} is sorted.
\end{exercise}

\begin{exercise}[$\star\star\star$]\label{exo:g11:proving-algorithms-correct:12}
Write out in full the proof that insertion sort sorts its list: the outer invariant, the inner invariant, why the
inner loop keeps it, and why both loops stop.
\end{exercise}

\begin{exercise}[$\star\star\star$]\label{exo:g11:proving-algorithms-correct:13}
Prove that the chapter's \code{search} returns an index of \code{x} whenever \code{x} is in the list.
\end{exercise}

\begin{exercise}[$\star\star\star$]\label{exo:g11:proving-algorithms-correct:14}
The loop \code{while n != 1}, with \code{n = n // 2} if \code{n} is even and \code{n = 3 * n + 1} otherwise, stops for
every \code{n} anyone has tried; from $27$ it takes $111$ turns. Why do neither \code{n} nor the number of turns
left serve as a \omterm{def:g11:proving-algorithms-correct:variant}{variant}? What would a proof that it always stops need?
\end{exercise}

\begin{exercise}[$\star\star\star$]\label{exo:g11:proving-algorithms-correct:15}
A function \code{f} gives \code{False} on \code{lo}, \code{lo + 1}, \dots{} up to some point, then \code{True}
up to \code{hi - 1}. Write \code{first_true(f, lo, hi)} that returns the first number where \code{f} is true, or \code{hi} if there is
none, by halving, with its invariant and its \omterm{def:g11:proving-algorithms-correct:variant}{variant}.
\end{exercise}

\section{Problem: The Integer Square Root}\label{sec:g11:proving-algorithms-correct:problem}

\begin{problem}[{Weekend problem --- an integer square root, proved and timed}]\label{pb:g11:proving-algorithms-correct:1}
The integer square root of a whole number $n$ is the largest whole number $r$ such that $r^2 \le n$: for $n = 15$ it is
$3$, for $n = 16$ it is $4$. The program below computes it twice, by counting up from $0$ and by halving a window, and
counts the turns of each.

\omcode{code/grade-11/05-proving-algorithms-correct/listings/isqrt.py}{1}{35}{Two integer square roots, checked against
Python's own on a million numbers.}\label{lst:g11:proving-algorithms-correct:isqrt}

\omoutput{code/grade-11/05-proving-algorithms-correct/outputs/isqrt.txt}{Each pair is the result and the number of
turns.}\label{out:g11:proving-algorithms-correct:isqrt}

\needspace{4\baselineskip}\medskip\noindent\textbf{Part I --- The specification.}
\begin{enumerate}
  \item Write the \omterm{def:g10:specifications-and-tests:pre-post}{precondition} and the \omterm{def:g10:specifications-and-tests:pre-post}{postcondition} of an integer square root function.
  \item Show that $r$ meets the \omterm{def:g10:specifications-and-tests:pre-post}{postcondition} exactly when $r^2 \le n < (r + 1)^2$.
  \item Give the integer square roots of $0$, $1$, $99$, $100$ and $101$.
  \item Why does such an $r$ exist, and why only one?
\end{enumerate}

\needspace{4\baselineskip}\medskip\noindent\textbf{Part II --- Counting up.}
\begin{enumerate}[resume]
  \item State an invariant of the loop of \code{isqrt_linear} that, at the end, gives the \omterm{def:g10:specifications-and-tests:pre-post}{postcondition}.
  \item Prove that it holds at the start and is kept.
  \item Show that \code{n - r * r} is a \omterm{def:g11:proving-algorithms-correct:variant}{variant}.
  \item How many turns does it take for $n$? For $n = 10^6$?
\end{enumerate}

\needspace{4\baselineskip}\medskip\noindent\textbf{Part III --- Halving.}
\begin{enumerate}[resume]
  \item Check that \code{lo * lo <= n < hi * hi} holds before the first turn.
  \item Show that each turn keeps it, in each of the two branches.
  \item Show that at the end \code{lo} meets the \omterm{def:g10:specifications-and-tests:pre-post}{postcondition}.
  \item Show that \code{mid} always lies strictly between \code{lo} and \code{hi} inside the loop.
  \item Deduce a \omterm{def:g11:proving-algorithms-correct:variant}{variant}, and that the loop stops.
\end{enumerate}

\needspace{4\baselineskip}\medskip\noindent\textbf{Part IV --- How many turns.}
\begin{enumerate}[resume]
  \item Follow the turns of \code{isqrt_halving(15)} and check the count printed.
  \item Show that if \code{hi - lo} is at most $2^k$ at the start of a turn, it is at most $2^{k - 1}$ at the end.
  \item Deduce that for $n + 1 \le 2^k$, the loop makes at most $k$ turns.
  \item Compute $2^{20}$; deduce a bound on the turns for $n \le 10^6$.
  \item Compare with the run: maximum and first $n$ reaching it. Why is the first such $n$ not $2^{19}$?
  \item What does the line \code{assert r == math.isqrt(n)} add to the proof, and what does the proof add to it?
  \item State the result: correctness, \omterm{def:g11:proving-algorithms-correct:variant}{termination} and the number of turns, against counting up.
\end{enumerate}
\end{problem}
